\documentclass[12pt]{article}
\usepackage[utf8]{inputenc}
\usepackage[T1]{fontenc}
\usepackage{hyperref}
\title{Waima'a Language Resources for K-12 Curriculum Development}
\author{Compiled by Claude Sonnet 5.5 (claude-sonnet-5-5) for Andr\'as Kornai}
\date{}
\begin{document}
\maketitle

\section{Introduction}
Waima'a (ISO 639-3 wmh; Glottocode waim1252) is spoken in Baucau Municipality, East Timor, between Vemasse and Bucoli and in and around Baucau city. Speaker estimates vary between sources (roughly 15,000 to 27,000). Its classification is described as uncertain by Omniglot: Malayo-Polynesian with much Papuan vocabulary, or a Papuan language related to Makasae. It is not known to be taught in schools. No full standalone dictionary or grammar was found in search results.

\section{Primary Educational Resources}
No primers, textbooks or school materials were found in search results.

\section{Linguistic and Reference Materials}
\begin{itemize}
\item \emph{Gl\'osariu Waima'a Caisido} (Belo, Bowden, Hajek, Himmelmann and Tilman; Baucau: Publica\c{c}\~oes Matebian, 2005). I saw a citation only, with no URL, so it is not in the JSON.
\item Hajek and Bowden (2010), on documenting and maintaining Waima'a, in \emph{Endangered Languages of Austronesia} (OUP). Citation only, no URL seen.
\item Omniglot overview: \url{https://www.omniglot.com/writing/waimoa.htm}
\item DoBeS poster on the Waima'a documentation: \url{https://dobes.mpi.nl/posters/print/oceania-waimaa-print.pdf}
\item DoBeS project page: \url{https://dobes.mpi.nl/?p=1381}
\item Endangered Languages Project entry: \url{https://www.endangeredlanguages.com/lang/2485}
\end{itemize}

\section{Religious and Community Resources}
None were found in search results.

\section{Digital and Audio Resources}
\begin{itemize}
\item PARADISEC collection KC2, which includes Waima'a and Midiki texts recorded in 2019 by Kirsten Culhane in Vemasse subdistrict: \url{https://admin-catalog.nabu-prod.paradisec.org.au/collections/KC2}
\item PARADISEC item wmh20191025: \url{https://admin-catalog.paradisec.org.au/collections/KC2/items/wmh20191025}
\item ``Waima'a Midiki recordings'' record: \url{https://researchdata.edu.au/waimaa-midiki-recordings/3976499}
\end{itemize}

\section{Curriculum Development Context}
The DoBeS project (2002--2006) worked with native speaker Maur\'icio Belo, who recorded, transcribed and translated into Tetum and Malay. The Caisido glossary and the archives above are the likeliest bases for materials. No school orthography or curriculum was found, and the speaker figures are uncertain.

\end{document}